% ID: SEL-D0080-OA-E02
% ESTADO: cotejada. Segunda pasada visual y borrador compilado revisados.
% Fuente DOC: Selectividad/sel_2012_matematicas_aplicadas/2 MATEMÁT. APLIC. CC. SS/EXAMEN 5-ANDALUCÍA.doc
% SHA-256 original: 87b729cdd4632386690a9390bdfa17dd7d5c902a2cf7e5801feea45e2b181969
% Página convertida: 1; opción A; ejercicio 2.

\begin{enumerate}[label=\alph*)]
\item \textbf{(0.75 puntos)} Para la función $f$ definida de la forma $f(x)=\dfrac{ax}{x+b}$, determine, razonadamente, los valores de $a$ y $b$ sabiendo que tiene como asíntota vertical la recta de ecuación $x=-2$ y como asíntota horizontal la de ecuación $y=3$.
\item \textbf{(1.75 puntos)} Para la función $g$, definida de la forma $g(x)=x^3-3x^2+2$, determine: su dominio, sus intervalos de crecimiento y decrecimiento y extremos relativos. Con esos datos haga un esbozo de su gráfica.
\end{enumerate}
